\begin{textsample}{POS Dim 1 – rn – Score 26.00 – rn/rn\_em\_17.txt}  \label{ex:f1_pos_239}
4.
JUVENTUDES, PROTAGONISMO E PROJETO DE VIDA NO ENSINO MÉDIO POTIGUAR “ Podemos afirmar \textbf{que} a juventude é plural, há diferentes modos de inserção dos jovens na vida social e cultural, portanto, é possível \textbf{dizer} \textbf{que} não há \textbf{uma} única juventude, mas sim, juventudes. ” ( DAYRELL \& GOMES, 2009 ; DAYRELL, 2003 ).
A etapa do Ensino Médio deve possibilitar aos jovens de diferentes origens e \textbf{que} \textbf{vivem} a juventude de modos díspares a troca de experiências, a \textbf{percepção} da vida social sob novas \textbf{óticas}, levando -os a refletirem, por exemplo, sobre questões étnico-raciais, de gênero, de trabalho e de sobrevivência, entre outras, em busca de sua identidade, tendo em vista \textbf{que} muitas vezes esses jovens vêm de percursos escolares tensos.
Esses \textbf{sujeitos} expõem nas escolas suas histórias de lutas por direitos sociais como trabalho, moradia, educação e \textbf{dignidade} de vida.
Absorvê -los em sua persistência de seguir \textbf{adiante} com sua formação escolar, alternando dedicação e esforços entre tempos de trabalho e de estudo, deve ser entendido como \textbf{uma} \textbf{vontade} de vencer, de obter conhecimento, de mudança de vida, de inclusão social, de busca por melhores oportunidades e conquista por direitos sociais e civis.
A escola deve abrir -se à compreensão das identidades e expressividades dos estudantes e das culturas juvenis \textbf{que} estes representam.
Consideremos nas juventudes do Rio Grande do Norte, suas identidades e expressividades, a cultura urbana, a cultura do grafite, dos jovens de periferia, do rock, do Hip Hop, do RAP, da literatura do cordel, da política estudantil, das juventudes engajadas nos movimentos religiosos ; do jovem do campo, \textbf{que} \textbf{consegue} conciliar o período de cultivo e os tempos de seca com as atividades escolares ; as dificuldades dos jovens oriundos e viventes nas comunidades quilombolas, suas lutas contra o preconceito e por políticas públicas \textbf{que} considerem suas especificidades ; a juventude indígena e a luta pela terra ; os ciganos e a sobrevivência de sua identidade, os jovens \textbf{que} estão na luta pelos direitos civis e inseridos na comunidade LGBTQIA+ e demais juventudes presentes no território potiguar.
Todas as juventudes deverão ser acolhidas, integralmente, pela escola e pelo Currículo.
No Ensino Médio, deparamo -nos com as várias expressividades das juventudes as quais se encontram e se reencontram, constroem e reconstroem atitudes, atos e relações \textbf{que} contribuirão para a formação cidadã desses \textbf{sujeitos}.
Frente a este contexto, “ a juventude ” é mais \textbf{que} \textbf{uma} palavra, \textbf{que} \textbf{um} período natural atribuído à idade, e, de fato, não deve ser entendida como singular, \textbf{precisa} e homogênea, mas como \textbf{um} marco social \textbf{historicamente} desenvolvido, a partir da multiplicidade de situações sociais, presente em \textbf{uma} etapa da vida, \textbf{que} condicionam as diferentes maneiras de “ ser jovem ”.
Segundo Margulis e Urresti ( 2008, p. 25 ) : > A condição de juventude manifesta -se de forma desigual conforme outros fatores, como classe social e/ou gênero.
Assim, não se devem considerar apenas os critérios biológicos de idade para definir juventude, não se pode também levar em conta apenas os critérios sociais ou apenas a dimensão simbólica.
Devem ser analisados também os aspectos fáticos, materiais, históricos e políticos em \textbf{que} toda produção social se desenvolve.
Essa essência \textbf{identitária} representa a concretude das múltiplas concepções de juventude, pautadas na perspectiva da diversidade, marcadas nas questões de gênero, nas diferentes regiões geográficas, dentre outras, concedendo a essa juventude \textbf{um} caráter social e culturalmente variável.
Destarte, ser jovem e estudante do EM \textbf{implica} sempre o entrecruzamento de papéis sociais diferenciados, vivenciados e apreendidos em outros lugares sociais, sem \textbf{que} ele despreze seu estilo particular de se expressar.
Esse fato gera desafios, como por exemplo, compreender as regras de convivência e lógica do ambiente escolar e, ao mesmo tempo, aprender a valorizar -se nesse espaço.
Este Referencial Curricular acolhe as juventudes do RN em sua multiplicidade e proporciona condições e oportunidades a todos os \textbf{sujeitos} dessa etapa de ensino, indistintamente, de cursar o Ensino Médio. \textbf{Um} exemplo dessas juventudes, são os estudantes oriundos de comunidades indígenas existentes no território potiguar.
Atualmente, o RN conta com 11 ( onze ) comunidades compostas por 7 ( sete ) povos integrantes do movimento indígena : Sagi-Trabanda, cujo povo se declara Potiguara ; Catu, onde se identificam Potiguaras Eleotérios ; Amarelão ; Serrote de São Bento ; Santa Terezinha, Marajó, Açuncena e Cachoeiras, comunidades formadas por \textbf{um} único povo \textbf{que} se denomina Mendonças Potiguaras do Amarelão ; Caboclos de Açú, \textbf{que} se afirmam índios Caboclos ; Lagoa do Tapará, formada pelos Tapuias da Lagoa do Tapará ou Tapuias Trarairiús ; e Lagoa do Apodi, cujos índios se reconhecem Tapuias Paiacús.³ ³ Conferir : MACEDO, A.
Índios do Rio Grande do Norte : vivos e em movimento.
Disponível em : https://www.saibamais.jor.br/indios-do-rio-grande-do-norte-vivos-e-em-movimento/.
Acesso em : 6 jun. 2020.
A SEEC compreende a pauta da educação indígena como \textbf{agenda} fundamental para a política educacional, a ser garantida às juventudes indígenas, no território do RN.
No \textbf{que} concerne à educação dos povos tradicionais, ela também é assegurada no Plano Nacional de Educação PNE, Lei nº 13.005/2014, quando estabelece em sua meta 9, estratégia para a oferta de educação para jovens e adultos, assim sendo, subentende -se também a população quilombola e aos povos indígenas.
A Constituição Brasileira, no § 1º, do Art. 215, determina \textbf{que} : > Art. 215.
O Estado garantirá a todos o pleno exercício dos direitos culturais e acesso às fontes da cultura nacional, e apoiará e incentivará a valorização e a difusão das manifestações culturais. > > § 1º O Estado protegerá as manifestações das culturas \textbf{populares}, indígenas e afro-brasileiras, e das de outros grupos participantes do processo civilizatório nacional. > > [... ] Em outro dispositivo constitucional, especificamente, o § 1º, II do Art. 225, a legislação materna garante a “ diversidade e a integridade do patrimônio genético do país ”, assim, os povos indígenas e comunidades quilombolas, grupos culturalmente diferenciados, devem receber políticas estatais \textbf{que} possibilitem condições sociais, culturais e econômicas na relação deles com o território e com o meio ambiente no qual estão inseridos.
Nesse sentido, o Referencial Curricular do território do RN intenciona garantir o direito do jovem aos conhecimentos de cada área, ao mesmo tempo em \textbf{que} reconhece haver \textbf{uma} permanente renovação de dados, informações, inovações, saberes por parte dos docentes e estudantes na condição de produtores de conhecimento.
As escolas e os currículos deixam de ser \textbf{meros} armazenadores de conteúdos informacionais, pelo contrário, são lugares de encontro de experiências sociais, reflexões, construção de leitura crítica \textbf{que} os estudantes fazem de si e do mundo e \textbf{que}, de forma \textbf{sistematizada}, buscam, de modo ativo, responder aos problemas da vida cotidiana \textbf{que} \textbf{necessitam} de respostas.
Assim, a SEEC, a partir do presente Referencial Curricular, a exemplo dos povos indígenas e quilombolas citados alhures, volve seu olhar para o acolhimento de todas as juventudes oriundas dos espaços sociais mais diversos e buscará o assentamento de políticas públicas educacionais \textbf{que} melhor atenda às juventudes existentes no nosso estado.
Para tanto, nossas escolas, do sertão ao litoral, precisam incentivar estes jovens a enfrentarem os desafios para construírem suas identidades e seus projetos de vida, promovendo espaços de mediação para atividades de lazer, práticas esportivas e produção de conhecimento, tornando estes espaços mais democráticos, \textbf{enquanto} lugares de interação, protagonismo, autonomia e resistências juvenis.
Percebemos, então, \textbf{que} é primordial \textbf{uma} boa interação entre professor e aluno pautada na exploração do conhecimento \textbf{sistematizado}, tecida pelos fios do aprender-ensinar-aprendendo.
Desse modo, a educação escolar fluirá melhor com o estudante exercendo o papel de parceiro e colaborador do professor.
Nessa perspectiva de reconhecimento e valorização do estudante do EM, a escola tem o papel fundamental de contribuir para \textbf{que} cada jovem se encontre e se perceba como protagonista de sua vida, respeite e se dê ao respeito em meio ao espaço de diversidade em \textbf{que} ele também se encontra e tome suas decisões filtradas em suas próprias crenças, conhecimentos e valores ; também é \textbf{tarefa} da escola acreditar e fazer o aluno acreditar, no desenvolvimento de seu potencial, buscando diferentes maneiras de se manter motivado ao desenho de seu projeto de vida.
Quando para \textbf{uns} o EM é \textbf{um} tempo de formação e sem grandes apreensões quanto à sua inserção no mercado de trabalho, para outros, essa relação estudo e trabalho apresenta -se conectada, tendo como motivo de preocupação sua projeção em \textbf{um} mercado cada vez mais competitivo e \textbf{que} tenha garantias profissionais atreladas a sua vida de estudante.
De acordo com o IBGE ( 2018 ), o mercado de trabalho norte-riograndense está marcado por significativas e persistentes desigualdades de gênero e \textbf{raça}, sendo este \textbf{um} aspecto \textbf{que} deve ser considerado nos processos de formulação, implementação e avaliação das políticas públicas, e, em particular, das políticas de emprego e inclusão social das juventudes do RN.
No seio das camadas \textbf{populares} é comum encontrarmos jovens \textbf{que} deixam de estudar e começam a trabalhar precocemente, muitos deles com idade inferior ao \textbf{que} determina a lei, ou, ainda, aqueles \textbf{que} estudam e trabalham ao mesmo tempo para poder ajudar no orçamento doméstico e manter suas famílias.
Para Corrochano, > A inserção precoce, a combinação entre trabalho e Ensino Médio e a postergação do ingresso no mercado de trabalho são resultados da interação de \textbf{um} conjunto de fatores.
Embora a necessidade de renda seja \textbf{um} fator \textbf{bastante} relevante para \textbf{que} muitos comecem a trabalhar antes da conclusão da escola média, outros aspectos também devem ser considerados, tais como : a conotação moral do trabalho, a conjuntura do mercado de trabalho, o sexo ( as chances de ser pressionado a trabalhar é maior entre rapazes ), a escolaridade dos pais, a ordem de nascimento na família, a quantidade de irmãos, o tipo de configuração familiar, a região de moradia, a experiência, dentre outros ( CORROCHANO, 2014, p.114 ).
Embora a escola não garanta o acesso ao universo dos empregos, em muitos casos ao ingressar no EM a \textbf{percepção} do jovem é a de \textbf{que} terá \textbf{um} caminho para o mercado de trabalho e, consequentemente, orientações para \textbf{um} futuro profissional.
Muitos ainda \textbf{veem} a escola como \textbf{um} local \textbf{que} poderia potencializar suas habilidades, interesses e descobertas para entrar em \textbf{um} mercado cada vez mais competitivo.
É \textbf{imprescindível} construir mecanismos de proteção aos jovens.
No âmbito geral, as escolas devem criar possibilidades na articulação com o mundo do trabalho e apoio na construção de projetos juvenis para o futuro profissional de todas as juventudes.
Assim, é necessário \textbf{que} a escola ofereça \textbf{um} ambiente com projetos e práticas pedagógicas com a finalidade de favorecer aos jovens no desenvolvimento de sua autonomia.
Preparar os estudantes para os desafios encontrados na sociedade \textbf{contemporânea}, está associado à importância \textbf{que} a escola atribui a eles \textbf{enquanto} \textbf{sujeitos} focais do processo educativo, pois, além de valorizar o seu olhar, cria diversas possibilidades para \textbf{que} se interessem pela resolução de problemas e contribuam para \textbf{uma} vida mais humanizada dentro da comunidade em \textbf{que} \textbf{vivem}.

% matched lemmas: adiante, agenda, bastante, conseguir, contemporâneo, dignidade, dizer, enquanto, historicamente, identitária, implicar, imprescindível, mero, necessitar, percepção, popular, preciso, que, raça, sistematizar, sujeito, tarefa, um, ver, viver, vontade, ótica
\end{textsample}
